% Part: computability
% Chapter: computability-theory
% Section: prop-reduce

\documentclass[../../../include/open-logic-section]{subfiles}

\begin{document}

\olfileid{cmp}{thy}{ppr}
\olsection{خواص قابلية الاختزال}

قولنا $A\leq_m B$ معناه أن $A$ تختزل إلى~$B$. وفي هذا الرمز إشارة:
متى كان $A\leq_m B$ لم تكن $A$ «أصعب من»~$B$، وكانت $B$ «لا تقل
صعوبة عن»~$A$. فكأن $\le_m$ ترتيب للمجموعات، أو لمسائل القرار، بحسب
تعقيدها، على مثال $\le$ في ترتيب الأعداد. والقضيتان الآتيتان تؤيدان هذا
المعنى؛ فأولاهما تثبت تعدي $\le_m$.

\begin{prop}
\ollabel{prop:trans-red}
إذا كان $A\leq_m B$ و$B\leq_m C$، فإن $A\leq_m C$.
\end{prop}

\begin{proof}
نركب اختزال $A$ إلى~$B$ مع اختزال $B$ إلى~$C$، فنحصل على اختزال $A$
إلى~$C$.
\end{proof}

\begin{prob}
أقم البرهان على \olref{prop:trans-red} بإثبات أنه متى كانت $f$ و~$g$ اختزالين
متعددين إلى واحد لـ$A$ إلى~$B$ ولـ$B$ إلى~$C$، على الترتيب، فإن
$\comp{f}{g}$ اختزال متعدد إلى واحد لـ$A$ إلى~$C$.
\end{prob}

\begin{prop}
\ollabel{prop:reduce}
إذا أخذنا مجموعتين $A$ و$B$، وفرضنا $A\leq_m B$، صح ما يأتي:
\begin{enumerate}
\item إذا كانت $B$ قابلة للتعداد حسابيًا، كانت~$A$ كذلك.
\item إذا كانت $B$ قابلة للحساب، كانت~$A$ كذلك.
\end{enumerate}
\end{prop}

\begin{proof}
لنأخذ $f$ اختزالًا متعددًا إلى واحد من $A$ إلى~$B$. أما الأمر
الأول، فلنجعل $B$ مجال دالة جزئية قابلة للحساب~$g$؛ فتكون $A$
مجال~$\comp{f}{g}$، لأن
\begin{align*}
x \in A & \text{ إذا وفقط إذا كان } f(x) \in B \\
& \text{ إذا وفقط إذا كانت }  g(f(x)) \fdefined.
\end{align*}

وأما الأمر الثاني، فإن قابلية $B$ للحساب تقتضي أن تكون $B$
و~$\Complement{B}$ قابلتين للتعداد حسابيًا
(\olref[cmp]{thm:ce-comp}). ومن اليسير أن نرى أن $f$ تصلح أيضًا
اختزال متعدد إلى واحد لـ$\Complement{A}$ إلى $\Complement{B}$؛ ولذلك
تكون $A$ و~$\Complement{A}$ !!{computably enumerable} بحسب الجزء الأول
من هذا البرهان. فتلزم قابلية $A$ للحساب من
\olref[cmp]{thm:ce-comp}. (ولنا سبيل آخر، وهو أن نتحقق من
$\Char{A}=\comp{f}{\Char{B}}$؛ فإذا كانت $\Char{B}$ قابلة للحساب،
كانت~$\Char{A}$ كذلك.)
\end{proof}

\begin{prob}
لتكن $f$ اختزالًا متعددًا إلى واحد لـ$A$ إلى~$B$. أثبت أن $f$ كذلك
اختزال متعدد إلى واحد لـ$\Complement{A}$ إلى $\Complement{B}$.
\end{prob}

\begin{prob}
أثبت، لكل اختزال $f$ متعدد إلى واحد من $A$ إلى~$B$، أن
$\Char{A}=\comp{f}{\Char{B}}$.
\end{prob}

\begin{digress}
ونحتاج في مباحث أخرى، وبخاصة براهين عدم القابلية للقرار، إلى مفهوم
أعم هو \emph{قابلية الاختزال بتورنغ}. فقد علمنا من
\olref[cmp]{cor:comp-k} لا تقبل متممة~$K_0$ الاختزال إلى~$K_0$، لأنها
غير قابلة للتعداد حسابيًا. غير أنك لو علمت أجوبة الأسئلة عن
$K_0$، لعرفت أجوبة الأسئلة عن متممتها أيضًا. ويقال إن مجموعة $A$
قابلة للاختزال بتورنغ إلى~$B$ إذا أمكن تعيين أجوبة الأسئلة عن~$A$
بإجراء قابل للحساب يستطيع طرح أسئلة عن~$B$. فهذا أوسع من
الاختزال متعددًا إلى واحد، الذي لا يجيز إلا (1) بطرح سؤال
واحد عن $B$، و(2) يجب أن تترجم إجابة «نعم» إلى إجابة «نعم» عن السؤال
المتعلق بـ$A$، وكذلك «لا». ويظل صحيحًا أنه إذا كانت $A$ قابلة للاختزال
بتورنغ إلى~$B$ وكانت $B$ قابلة للحساب، فإن $A$ قابلة للحساب أيضًا
(مع أن العبارة المناظرة لا تصح، كما رأينا، في القابلية للتعداد حسابيًا).

فتأمل وجوه الاختزال المختلفة، واستبن ما بينها من
فروق. وإنما نقصر هذا الفصل على الاختزال
متعددًا إلى واحد. وللنوعين المذكورين في الفقرة
السابقة نظيران في التعقيد الحسابي إذا زيد شرط عمل
آلات تورنغ في زمن كثير الحدود: فنسخة التعقيد من قابلية الاختزال متعددًا
إلى واحد تعرف بـ\emph{قابلية الاختزال بكارب}، أما نسخة التعقيد من
قابلية الاختزال بتورنغ فتعرف بـ\emph{قابلية الاختزال بكوك}.
\end{digress}

\end{document}
